% =============================================================================
% 7. LIMITATIONS AND DISCUSSION
% =============================================================================
\section{Limitations and Discussion}
\vspace{-0.5ex}
\label{sec:limitations}
We close by stating four limitations that scope our claims.
(i) \textit{Python-only corpus and evaluation.} The mid-training corpus
and the in-domain agent benchmarks are exclusively Python; cross-language
evidence comes only indirectly through FullStackBench-EN
(Section~\ref{sec:exp:gen}), and transfer to Java, C++, or Rust is left
to future work.
(ii) \textit{Teacher dependency for CoT.} The default recipe relies on
Gemini-3-Flash; the CoT-source ablation (Section~\ref{sec:exp:ablation})
shows self-generated rationales recover most of the gain, but a fully
open-source replication requires a comparably strong open teacher.
(iii) \textit{Partial cross-base validation.} Our cross-base evidence
comes from a single non-Qwen2.5-Coder configuration (Qwen3-8B with
SWE-Lego), which simultaneously varies the post-training pipeline; the
result indicates the recipe is not tied to one pretraining/post-training
combination rather than guaranteeing transfer across all base families.
(iv) \textit{Modularity assumption.} Function-aware FIM presupposes
modular code; on monolithic scripts, generated code, or notebooks, the
pipeline yields fewer eligible targets, a regime we do not study
systematically.
